\documentclass[10pt,letterpaper,twocolumn]{article}
\usepackage[margin=0.72in,top=0.70in,bottom=0.72in]{geometry}
\usepackage[T1]{fontenc}
\usepackage{newtxtext,newtxmath}
\usepackage{microtype}
\usepackage{titlesec}
\usepackage{enumitem}
\usepackage[hidelinks]{hyperref}
\usepackage{fancyhdr}
\usepackage{lastpage}
\usepackage{comment}


\setlength{\columnsep}{0.25in}
\setlength{\parindent}{0.18in}
\setlength{\parskip}{0pt}

\titleformat{\section}{\normalfont\scshape\centering}{\thesection.}{0.45em}{}
\titlespacing*{\section}{0pt}{0.8em}{0.42em}

\renewcommand{\thesection}{\Roman{section}}

\pagestyle{fancy}

\fancyhf{}
\fancyhead[R]{\itshape RIT Computer Science \enspace\textbullet\enspace Capstone Report}
\fancyfoot[C]{Rochester Institute of Technology}
\fancyfoot[R]{\thepage\space|\space Page}

\renewcommand{\headrulewidth}{0pt}
\renewcommand{\footrulewidth}{0.2pt}

\setlength{\headheight}{14pt}

\begin{document}

\title{\LARGE\bfseries Distributed Training of Large Language Models\\[0.10em]\par}

\author{
    \begin{minipage}[t]{0.45\textwidth}
        \centering
        Andrew Photinakis \\
        Rochester Institute of Technology \\
        \texttt{acp7795@rit.edu} \\
        Rochester, New York, USA
    \end{minipage}
}

\maketitle

\section{Abstract}
\begin{comment}
### 1. Title & Abstract
* **Title**: **10 to 15 words** long. It should be concise and descriptive of your specific focus area (e.g., *“Distributed Training of Large Language Models: Architectures, Parallelism, and Communication Challenges”*).
* **Abstract**: **6 to 8 lines** long. It must be a self-contained summary written so that a computer science layperson can understand the core problem, context, and key conclusions of the paper.
\end{comment}

\section{Introduction}

\section{Distributed Architectures \& Algorithms}

\section{Research Advances \& Synthesis}

\section{Scope, Critique \& Challenges}

\section{Conclusion}

\section{References}

\end{document}
